% mazzi: 33
\chapter{La storia della preparazione fisica nel calcio}
\label{cap:storiaprepfisica}

Il percorso della preparazione fisica nel calcio è la storia di un passaggio: da una fase
\textbf{empirica e artigianale}, in cui l'allenamento fisico era affidato all'esperienza personale
degli ex-calciatori diventati allenatori, a una fase \textbf{scientifica e ``industriale''}, in cui
metodiche, strumenti e carichi sono ricavati dalla conoscenza fisiologica del giocatore e dallo
studio dell'entità dello sforzo in gara.

% ---------------------------------------------------------------------------
\section{Le origini del gioco}

\rubrica{I precedenti più remoti}
\begin{itemize}
  \item \textbf{kemari} (antico Giappone) e \textbf{tsu-chu} (antica Cina): i precedenti più remoti
    del gioco del calcio;
  \item datazione incerta: le tradizioni locali parlano di un migliaio di anni prima di Cristo,
    altre fonti collocano il tsu-chu molto più indietro, attorno al \textbf{2600 a.C.};
  \item elementi comuni ai due giochi:
    \begin{itemize}
      \item uso dei \textbf{piedi};
      \item presenza di una \textbf{``porta''} rudimentale, definita da due alberi o da aste di
        bambù;
      \item utilizzo di una \textbf{palla}.
    \end{itemize}
  \item \textit{chu}: palla di cuoio realizzata con la vescica di animale gonfiata, oppure riempita
    di capelli femminili;
  \item nel \textbf{V secolo a.C.} il tsu-chu faceva parte dei programmi di \textbf{addestramento
    militare} dell'esercito ed era quindi finalizzato --- come molti altri esercizi ---
    all'\textbf{efficienza fisica dei soldati}.
\end{itemize}

\rubrica{La nascita del calcio moderno}
\begin{itemize}
  \item \textbf{26 ottobre 1863, Inghilterra}: nasce la \textbf{Football Association}, e con essa
    ufficialmente il calcio moderno;
  \item l'elaborazione plurisecolare del gioco si fissa in un atto ufficiale: \textbf{undici
    dirigenti} di club e scuole londinesi, riuniti nella \textit{Free Masons Tavern} sulla
    \textit{Great Queen Street}, fondano la Football Association;
  \item \textbf{1897}: a Londra viene istituita la prima \textbf{associazione di giocatori
    britannici}, poi trasformatasi nella potente \textbf{PFA} (\textit{Professional footballer's
    association}) --- un passo importante verso il professionismo.
\end{itemize}

\rubrica{La FIFA e le federazioni continentali}
\begin{itemize}
  \item \textbf{1904, Parigi}: nasce la \textbf{FIFA} (\textit{Federation Internazionale de Football
    Association}) grazie ai rappresentanti di \textbf{sette Associazioni nazionali} --- Francia,
    Belgio, Olanda, Svizzera, Danimarca, Svezia e Spagna;
  \item obiettivo: \textbf{rendere unico il calcio} attraverso lo stesso regolamento. La FIFA
    diventa l'\textbf{unico ente in grado di modificare le regole di gioco}, dando notevole
    credibilità e impulso alla crescita del calcio;
  \item conseguenza immediata: diventa possibile organizzare \textbf{partite tra squadre e
    rappresentative di Nazioni diverse};
  \item oggi la FIFA ha in ogni continente un'\textbf{appendice} che regola i campionati
    continentali per Nazioni e per club: per l'Europa è l'\textbf{UEFA}, con sede a \textbf{Nyon}
    (Svizzera).
\end{itemize}

% ---------------------------------------------------------------------------
\section{Dalle origini agli anni Cinquanta}

\subsection{I primi riferimenti alla resistenza}

In Europa, solo \textbf{dopo la Prima guerra mondiale} si cominciò a considerare la necessità di una
buona preparazione fisica ai fini prestativi del calciatore.

\begin{itemize}
  \item \textbf{Szehany} (ungherese), \textit{Tecnica del giuoco del calcio} (\textbf{1928}):
    \begin{itemize}
      \item mette in rilievo l'importanza della \textbf{resistenza};
      \item consiglia \textbf{corsa di durata}, \textbf{nuoto}, \textbf{ciclismo};
      \item si tratta di consigli \textbf{senza alcun indirizzo sistematico}, ma la resistenza
        viene considerata il fattore determinante della prestazione calcistica.
    \end{itemize}
  \item \textbf{Ermanno Morelli} (\textbf{1934}):
    \begin{itemize}
      \item critica gli allenatori delle grandi società, cui imputa l'\textbf{empirismo degli
        ex-calciatori};
      \item afferma che occorre \textbf{variare il lavoro} in base al \textbf{ruolo}, alle
        \textbf{caratteristiche fisiche} e al \textbf{peso corporeo};
      \item per i giocatori in sovrappeso consiglia \textbf{salti alla corda per un'ora intera},
        con intervalli di 5 o 10 minuti.
    \end{itemize}
\end{itemize}

\subsection{La differenziazione per ruolo nei ritiri della nazionale}

Nei ritiri della nazionale italiana --- campione del mondo \textbf{1934} e \textbf{1938}, alloro
olimpico a \textbf{Berlino 1936} --- la preparazione fisica, o \textbf{attività ginnico-atletica}
come si era soliti definirla, veniva effettuata \textbf{a gruppi} e variava in base ai ruoli:

\begin{itemize}
  \item \textbf{portieri}: esercizi di \textbf{agilità} e salti alla corda, \textbf{senza} giri di
    campo né attività di corsa --- per il ruolo del portiere la resistenza non era considerata
    importante;
  \item \textbf{terzini}: brevi tratti di \textbf{corsa veloce}, esercitazioni di salto, ginnastica
    di mobilizzazione;
  \item \textbf{mediani}: \textbf{corsa di fondo} alternata con la marcia, fino a un totale di
    \textbf{4--5 km} per seduta, più alcune ripetizioni di sprint;
  \item \textbf{attaccanti}: massima importanza agli \textbf{esercizi di scatto}, con molteplici
    varianti (arresti repentini, cambi di direzione).
\end{itemize}

\subsection{La stasi fino alla Seconda guerra mondiale}

\begin{itemize}
  \item sul tema dell'allenamento fisico non sembra si siano tentate \textbf{nuove strade} né
    sperimentati \textbf{nuovi metodi};
  \item il maggiore interesse calcistico era orientato sulla \textbf{dialettica dei sistemi di
    gioco}: era la \textbf{tattica} ad avere la predominanza sull'aspetto fisico;
  \item proprio l'introduzione di \textbf{schemi più razionali} --- nella distribuzione dei
    giocatori sul campo e nei ruoli loro assegnati --- aveva però cominciato a sollecitare il
    problema di un \textbf{allenamento fisico specifico} per le forme di movimento che il
    particolare compito richiedeva.
\end{itemize}

\subsection{Il dopoguerra: la preparazione atletica diventa essenziale}

Subito dopo la Seconda guerra mondiale la preparazione atletica comincia a essere considerata una
\textbf{componente essenziale} dell'allenamento calcistico, e l'attenzione degli esperti si orienta
sulle metodiche più adatte a portare il calciatore, nel modo più rapido e continuativo possibile, a
esprimere la sua \textbf{``vitalità di gioco''}:

\begin{itemize}
  \item arrivare \textbf{per primi sulla palla};
  \item \textbf{saltare più in alto} dell'avversario;
  \item \textbf{resistere alle cariche} dell'avversario;
  \item mantenere \textbf{inalterati i riflessi} per tutta la durata dell'incontro.
\end{itemize}

\rubrica{Platko: il primo approccio sistematico}
\begin{itemize}
  \item \textbf{Fransisco Platko} (ungherese), \textit{Arte y ciencia del futbol moderno} (Santiago
    del Cile, \textbf{1946}): affronta per la prima volta il problema della preparazione atletica
    nel calcio in modo \textbf{sistematico e razionale};
  \item parte dalla constatazione che ``è un problema molto discusso se il giocatore avesse bisogno
    di poco o molto allenamento'';
  \item con notevole intuito insiste sul fatto che l'allenamento debba essere caratterizzato
    dall'\textbf{alternanza di fasi intense con altre relativamente leggere};
  \item per l'attività di corsa consiglia di \textbf{alternare corse rapide con corse lente}.
\end{itemize}

\rubrica{Cornilli: il modello di seduta}
\textbf{Jean Cornilli}, istruttore ai corsi per allenatori in Francia, in \textit{Le football
dévoilé} (\textbf{1951}) descrive un modello di seduta di allenamento:
\begin{enumerate}
  \item messa in azione;
  \item esercizi addominali;
  \item esercizi di agilità e di acrobatica elementare (capovolte, ecc.);
  \item esercizi tecnici con la palla;
  \item defaticamento.
\end{enumerate}

\rubrica{Le pubblicazioni FIGC (1950--1952)}
Si attengono allo stesso metodo e, corredate da nozioni generali sui criteri di valutazione, sulla
\textbf{morfologia} e sulla \textbf{costituzione umana}, mettono in evidenza le caratteristiche
generali della \textbf{contrazione muscolare statica e dinamica}. Il gioco del calcio comincia a
essere \textbf{studiato e analizzato nelle sue componenti peculiari}:
\begin{itemize}
  \item esercizi di \textbf{forza} alternati con altri di \textbf{allungamento} e
    \textbf{mobilità articolare};
  \item si introduce il concetto di \textbf{preatletismo generale}, orientato al potenziamento
    delle \textbf{qualità muscolari} e di quelle \textbf{organiche};
  \item viene presa in esame la \textbf{tecnica della corsa e del salto}, con particolare
    riferimento ai movimenti e alle esigenze calcistiche.
\end{itemize}

% ---------------------------------------------------------------------------
\section{Dal 1954 ai Mondiali del 1970}

\subsection{1954: la vittoria della condizione fisica sulla tattica}

\begin{itemize}
  \item \textbf{Campionati del Mondo 1954} (Svizzera): vince la \textbf{Germania Occidentale} di
    \textbf{Sepp Herberger}, che nella finale contro l'Ungheria lascia stupefatti gli spettatori di
    tutto il mondo (per la prima volta le partite vengono trasmesse in televisione);
  \item il successo è universalmente considerato una \textbf{vittoria della condizione fisica sulla
    tattica};
  \item contrariamente a quanto sarebbe stato logico attendersi --- cioè l'esaltazione del
    condizionamento fisico dal punto di vista \textbf{organico} (cioè della resistenza) e
    \textbf{muscolare} --- l'idea prevalente
    resta però che la forma di allenamento più idonea consista nel \textbf{ripetere metodicamente i
    movimenti} che il calciatore sarebbe stato chiamato a compiere durante la gara.
\end{itemize}

\subsection{Weisweiler: l'allenamento specifico attraverso il gioco}

\textbf{Hennes Weisweiler} (allenatore del Borussia Mönchengladbach), nel capitolo dedicato alla
condizione fisica di \textit{Der Fussbal-Taktik, Training und Mannschaft} (\textbf{1959}), si
dichiara \textbf{contrario} all'utilizzo, nell'allenamento dei calciatori, delle metodiche
specifiche dell'\textbf{atletica leggera}.

\begin{itemize}
  \item il calciatore \textbf{non è} né un velocista, né un mezzofondista, né un maratoneta, e
    nemmeno un sollevatore di pesi: \textbf{deve semplicemente giocare a calcio};
  \item gli occorrono \textbf{forza}, \textbf{velocità} e \textbf{resistenza} necessarie a
    utilizzare la propria \textbf{tecnica al servizio della squadra} $\rightarrow$ serve un
    \textbf{allenamento specifico}, che non può essere preso a prestito altrove;
  \item riconosce il \textbf{valore accessorio} di un buon addestramento alla corsa e al salto e di
    una adeguata ginnastica di rafforzamento per i muscoli del \textbf{tronco}, delle
    \textbf{spalle} e delle \textbf{gambe};
  \item insiste sulla \textbf{preponderante efficacia delle forme di allenamento basate sui
    giochi}, con e senza palla: è soltanto per mezzo del gioco che il calciatore può essere
    stimolato --- anche dal punto di vista \textbf{psicologico} --- a partecipare a un allenamento
    sufficientemente \textbf{intenso} e opportunamente \textbf{vario};
  \item attività calcistiche fondamentali: le esercitazioni (giochi) \textbf{3:1}, \textbf{4:2},
    \textbf{3:2}, \textbf{1:1}, \textbf{3:3}, oltre a forme ridotte e semplificate di
    \textbf{pallamano} e \textbf{pallacanestro}.
\end{itemize}

\subsection{1958: il Brasile e la disciplina atletica}

\begin{itemize}
  \item i \textbf{Mondiali del 1958} (Svezia) sono vinti dal \textbf{Brasile} grazie all'eccelso
    \textbf{valore tecnico} dei giocatori: i due Santos, Didì, Vavà, Pelè, Garincha, con qualità
    tecnica e fantasia da non temere rivali in tutto il mondo;
  \item in uno studio sulla preparazione delle squadre partecipanti, l'ungherese \textbf{Arpad
    Csamadi} afferma di essere rimasto \textbf{sorpreso dalla serietà e dalla disciplina} con cui i
    giocatori brasiliani --- considerati generalmente poco disposti a forme di allenamento intense
    e impegnative --- seguivano le esercitazioni anche di natura squisitamente atletica.
\end{itemize}

\subsection{I corsi internazionali UEFA: circuito e interval training}

\rubrica{1961, Macolin (Svizzera): 1\textsuperscript{o} corso internazionale per allenatori}
Organizzato dalla Commissione tecnica dell'UEFA.
\begin{itemize}
  \item \textbf{Walter Winterbottom} (per anni CT dell'Inghilterra) illustra e dimostra
    praticamente due forme di allenamento basate sul \textbf{``metodo del circuito''}:
    \begin{itemize}
      \item una per la \textbf{resistenza muscolare};
      \item una per la \textbf{resistenza organica}.
    \end{itemize}
  \item molti allenatori cominciano a metterlo in pratica sul campo, inserendolo nei propri
    programmi;
  \item limite: il circuito doveva comprendere \textbf{tutta} l'attività inerente la condizione
    fisica, \textbf{sostituendo completamente ogni altra forma di esercitazione}.
\end{itemize}

\rubrica{1962, Hennef-Sieg Bonn: 2\textsuperscript{o} corso internazionale}
\begin{itemize}
  \item alla Scuola dello sport, sempre Winterbottom presenta un programma di \textbf{interval
    training}, il cosiddetto \textbf{``dei 45 secondi''};
  \item è uno fra i primi tentativi di applicare all'allenamento calcistico una metodica che a quel
    tempo godeva di grande prestigio, per i risultati già conseguiti in \textbf{altre discipline
    sportive};
  \item \textbf{non trova largo seguito}, anche per alcune \textbf{deviazioni dalla forma
    originaria} che lo rendevano intenso e faticoso.
\end{itemize}

\subsection{1966--1967: il calcio atletico}

\begin{itemize}
  \item \textbf{1966, Londra}: la finale del Campionato del Mondo vede contrapposte
    \textbf{Inghilterra} e \textbf{Germania Occidentale}, le due nazioni che già da alcuni anni
    venivano giudicate le tipiche rappresentanti del cosiddetto \textbf{calcio atletico};
  \item i Mondiali d'Inghilterra evidenziano che il calcio moderno non può più basarsi solo su una
    \textbf{tecnica raffinata} o su accorgimenti \textbf{puramente tattici}, ma richiede anche e
    soprattutto giocatori dotati di \textbf{elevato dinamismo};
  \item \textbf{1967, convegno UEFA} (Zeist, Olanda): i relatori più seguiti sono l'inglese
    \textbf{Allen Wade} e il tedesco \textbf{Helmut Schon}. Le loro relazioni dimostrano
    l'importanza della preparazione atletica nel calcio e dello sviluppo, attraverso l'esercizio,
    delle qualità fisiche richieste al calciatore:
    \begin{itemize}
      \item resistenza muscolare;
      \item resistenza organica;
      \item forza;
      \item velocità.
    \end{itemize}
  \item lo scozzese \textbf{Roy Small} presenta una seduta di \textbf{interval training calcistico}:
    \begin{itemize}
      \item distanza di \textbf{100 m}, ripetuta \textbf{15 volte};
      \item intervallo di recupero variabile fra \textbf{60 e 75 s} fra una ripetizione e l'altra;
      \item il programma completo prevedeva l'inserimento di questa tipologia di lavoro per
        \textbf{quattro settimane consecutive} durante il periodo \textbf{precampionato}.
    \end{itemize}
\end{itemize}

\subsection{1970, Messico: la ricerca entra nei Mondiali}

\begin{itemize}
  \item i Campionati del Mondo del \textbf{1970} in Messico richiamano al seguito di atleti e
    squadre \textbf{medici, biologi, ricercatori};
  \item obiettivo principale: esaminare le \textbf{reazioni dell'organismo umano in condizioni di
    massima prestazione in altitudine};
  \item veri e propri \textbf{laboratori scientifici} vengono organizzati in loco, per svolgere le
    ricerche direttamente alle elevate quote dell'altipiano messicano;
  \item la nazionale italiana ottiene il titolo platonico di \textbf{vicecampione del Mondo}
    (battuta in finale dal Brasile 4--1), dopo aver superato la Germania Occidentale
    \textbf{4--3} al termine dell'epica partita protrattasi per oltre \textbf{due ore}.
\end{itemize}

% ---------------------------------------------------------------------------
\section{Dai Mondiali del Messico a oggi}

\subsection{Il calcio totale e la nascita della figura del preparatore}

\begin{itemize}
  \item negli ultimi trent'anni si sono verificati notevoli cambiamenti sul piano
    \textbf{tecnico-tattico} e \textbf{fisico};
  \item ciò è avvenuto grazie al \textbf{calcio totale}:
    \begin{itemize}
      \item praticato dagli \textbf{Olandesi} ai Mondiali del \textbf{1974} in Germania Ovest;
      \item e dal \textbf{Milan di Sacchi} nel periodo \textbf{1987--91}.
    \end{itemize}
  \item concetti di gioco affermati in maniera decisiva:
    \begin{itemize}
      \item il \textbf{pressing};
      \item i \textbf{raddoppi di marcatura};
      \item la ricerca costante della \textbf{superiorità numerica}, in fase sia difensiva sia
        offensiva, tramite il \textbf{movimento continuo} di tutti i giocatori della squadra.
    \end{itemize}
  \item lo studio della preparazione fisica acquisisce sempre maggiore importanza e diventa sempre
    più \textbf{accurato e scientifico}, specialmente in Italia dove, dal \textbf{1991}, il
    \textbf{Settore Tecnico della FIGC} organizza annualmente a \textbf{Coverciano} un corso
    specifico per l'\textbf{abilitazione a preparatore atletico} del calcio;
  \item la figura del preparatore è ormai quasi indispensabile nel difficile compito di:
    \begin{itemize}
      \item far raggiungere ai calciatori un buon livello di \textbf{forma già all'inizio del
        campionato};
      \item e, ciò che più conta, far \textbf{mantenere uno standard di performance} per i mesi
        successivi di intensa attività agonistica.
    \end{itemize}
\end{itemize}

\subsection{La compressione del precampionato}

\begin{itemize}
  \item cambiamento importante degli ultimi anni: la crescente importanza di \textbf{tornei} o
    \textbf{tournée} che si svolgono già dopo \textbf{pochi giorni} di preparazione precampionato;
  \item conseguenza: per diverse squadre il periodo da dedicare alla \textbf{sola preparazione} è
    assai più limitato rispetto al passato, quando in alcune settimane di training erano inserite
    soltanto partite di \textbf{importanza minima o nulla} --- le amichevoli fra la squadra A e la
    squadra B, o contro squadre di categoria inferiore della zona in cui si svolgeva il ritiro;
  \item per ovviare alla carenza di tempo da dedicare allo sviluppo delle qualità fisiche nel
    precampionato è diventato sempre più frequente (nei professionisti):
    \begin{itemize}
      \item essere seguiti \textbf{anche in vacanza} da un preparatore della società;
      \item o addirittura svolgere veri e propri \textbf{pre-ritiri}.
    \end{itemize}
\end{itemize}

% ---------------------------------------------------------------------------
\section{La strumentazione per la valutazione e il controllo dell'allenamento}

Con l'inserimento a pieno titolo dei preparatori nell'organico delle squadre di calcio si è
intensificata la \textbf{sperimentazione sulla metodologia dell'allenamento}, con il supporto sia
delle \textbf{teorie sviluppate in laboratorio} sia della \textbf{ricerca tecnologica}, cui si deve
la diffusione di strumenti di grande efficacia per la \textbf{valutazione funzionale} dell'atleta e
per il \textbf{controllo dell'allenamento}. È diventato abituale l'impiego di:

\begin{itemize}
  \item \textbf{cardiofrequenzimetri}: determinano con precisione lo sforzo e quindi il
    \textbf{carico di lavoro} di una determinata esercitazione;
  \item \textbf{GPS}: valutano e quantificano il carico proposto ai giocatori, misurando la
    \textbf{velocità di spostamento}, rilevata da \textbf{10, 15 o 20 volte al secondo} a seconda
    della \textbf{frequenza di campionamento} --- 10, 15 o 20 Hz, cioè altrettanti punti misurati
    al secondo;
  \item \textbf{lattacidometri}: rilevazione immediata della quantità di \textbf{acido lattico}
    (lattato) presente nel sangue post esercitazione;
  \item \textbf{plicometri}: calcolo della \textbf{percentuale di massa grassa} dell'atleta;
  \item \textbf{pedana computerizzata di Bosco}: determina i valori delle varie \textbf{espressioni
    di forza};
  \item \textbf{Muscle Lab}: testando l'atleta con una \textbf{serie di carichi crescenti} fornisce
    il grafico della relazione \textbf{forza / velocità / potenza}, da cui si ricavano:
    \begin{itemize}
      \item la \textbf{1 RM} (carico massimo);
      \item il \textbf{picco di potenza} per ogni carico;
      \item il \textbf{rapporto ottimale fra forza e velocità}.
    \end{itemize}
  \item \textbf{fotocellule in telemetria}: verificano i \textbf{tempi esecutivi} delle varie
    tipologie di esercitazione --- accelerazioni sui \textbf{5--10 m}, sprint sui \textbf{20--30
    m}, navette, ecc.
\end{itemize}

% ---------------------------------------------------------------------------
\section{Dall'empirico allo scientifico: l'inversione della piramide}
\label{sec:inversione-piramide}

\rubrica{La trasformazione}
Il mondo del calcio, per quanto concerne l'allenamento delle qualità fisiche, ha registrato una
profonda trasformazione, passando da una fase \textbf{empirica e artigianale} a una che si può
definire \textbf{scientifica e ``industriale''}, cui si devono i notevoli passi avanti riscontrabili
nella preparazione fisica.

Tale trasformazione ha determinato una \textbf{razionalizzazione dell'allenamento}, con la messa a
punto di:
\begin{itemize}
  \item \textbf{metodiche più efficaci} per lo sviluppo delle qualità fisiche utili al calciatore,
    fondate su una conoscenza più precisa delle \textbf{caratteristiche fisiologiche} e
    dell'\textbf{entità dello sforzo del giocatore} in gara e in allenamento;
  \item l'applicazione sui \textbf{singoli atleti} di carichi e tipi di lavoro adatti alle loro
    \textbf{caratteristiche individuali} e alle loro \textbf{specifiche carenze e necessità}.
\end{itemize}

\rubrica{Dalla resistenza alla forza}
\begin{itemize}
  \item in passato --- anche relativamente recente --- la metodologia dell'allenamento era
    orientata prevalentemente allo sviluppo della \textbf{resistenza};
  \item attualmente, sia in precampionato sia durante la fase agonistica, il \textbf{potenziamento
    muscolare} --- soprattutto dei gruppi muscolari maggiormente sollecitati nella prestazione
    calcistica --- è diventato uno degli argomenti principali nelle programmazioni dei preparatori;
  \item è ormai opinione ampiamente diffusa che l'allenamento debba essere basato \textbf{più sulla
    forza che sulla resistenza}, perché i momenti più significativi di una gara sono quelli in cui
    il calciatore compie prestazioni influenzate dalla forza:
    \begin{itemize}
      \item tiri;
      \item scatti, accelerazioni e decelerazioni;
      \item arresti e cambi di direzione improvvisi;
      \item contrasti;
      \item stacchi per colpire di testa.
    \end{itemize}
\end{itemize}

\rubrica{L'inversione della piramide}
La conclusione del percorso storico si riassume in uno schema a \textbf{piramide} su tre livelli,
i cui estremi sono le \textbf{fibre rosse} (base) e le \textbf{fibre bianche} (vertice). La
piramide si può percorrere in due direzioni opposte.

\begin{itemize}
  \item i tre livelli, dal basso verso l'alto:
    \begin{enumerate}
      \item \textbf{base} --- corse condotte sotto controllo della \textbf{frequenza cardiaca},
        attorno all'\textbf{80\% della FC massima}: è il livello delle \textbf{fibre rosse};
      \item \textbf{livello intermedio} --- esercitazioni di corsa \textbf{senza palla};
      \item \textbf{vertice} --- lavoro \textbf{con la palla}, cioè l'attività specifica del
        calciatore: è il livello delle \textbf{fibre bianche}.
    \end{enumerate}
  \item \textbf{percorso tradizionale} (dal basso verso l'alto), praticato fino a qualche decennio
    fa:
    \begin{itemize}
      \item si comincia la preparazione con \textbf{corse lunghe e continue}, a ritmo costante e
        con un \textbf{volume di lavoro sostenuto};
      \item obiettivo di partenza: l'allenamento delle \textbf{fibre rosse};
      \item solo in un secondo momento si sale verso il lavoro specifico.
    \end{itemize}
  \item \textbf{inversione della piramide} (dall'alto verso il basso), impostazione attuale:
    \begin{itemize}
      \item si comincia \textbf{fin da subito con il lavoro con la palla}, cioè con l'attività
        specifica del calciatore;
      \item la \textbf{resistenza} non è più un aspetto prioritario e fondamentale, ma
        \textbf{secondario};
      \item la si sviluppa quindi per via \textbf{indiretta} più che diretta.
    \end{itemize}
\end{itemize}
